\addchap{Abstract}

Emergency response depends on decisions made quickly with incomplete information. In Egypt, a
person in distress must choose between separate numbers for the police, the ambulance service
and the fire brigade; dispatchers work from spoken descriptions that are hard to prioritise; and
crews have little way to prepare the hospital that will receive a patient. These gaps cost
minutes that matter.

This book presents \textbf{\ProjectName} (\ProjectFullName), a complete emergency-response
platform that addresses them. A citizen files a report with one tap; the platform captures
the location, classifies the report's priority and flags possible duplicates with a
language model, and places it in a live queue for dispatchers. A dispatcher offers the
work to the nearest suitable response unit; the crew accepts, travels, arrives and completes
the mission from a mobile app, while the reporter follows every step live and can talk to the
responders. Ambulance crews choose a receiving hospital and stream the patient's vital signs
to it during transport. The platform serves ten user roles through a Flutter mobile app and a
Next.js web console, backed by a Spring Boot API with a PostgreSQL database, a separate AI
model service and managed cloud services.

The design concentrates on the qualities an emergency service cannot do without. Report
submission is idempotent and the mobile app queues reports taken offline; notifications are
emitted only after a transaction commits; a graduated caller-trust policy slows repeat abusers
but never refuses a genuine emergency; and a patient's medical profile is shared only with
consent, only with people treating the patient, only while the incident is open, and with every
read audited. Sensitive content is encrypted at rest, and the system works in English and
Arabic with right-to-left layout.

Verification combines automated unit tests of the safety-critical rule classes, continuous
integration, and a catalogue of end-to-end scenarios traced to the requirements. The work shows
that a small team can build, with modern engineering practice and managed cloud services, a
coherent platform that connects citizens, dispatchers, field crews and hospitals in a single,
auditable workflow.

\vspace{0.3cm}
\noindent\textbf{Keywords:} emergency dispatch; decision support; real-time systems;
large language models; mobile computing; privacy by design; software engineering.
\clearpage
